\documentclass[12pt,a4paper]{article}
\usepackage[utf8]{inputenc}
\usepackage[T1]{fontenc}
\usepackage{amsmath}
\usepackage{amssymb}
\usepackage{graphicx}
\usepackage[spanish]{babel}
\usepackage[a4paper, margin=2.5cm]{geometry}
\usepackage[european]{circuitikz}
\usetikzlibrary{babel}

\title{Análisis de Equivalentes de Thévenin y Norton}
\author{Gemini Code Assist and Prof. CERO}
\date{\today}

\begin{document}

\maketitle

\section{Descripción del Circuito}

El objetivo es determinar el circuito equivalente de Thévenin y de Norton entre los terminales 
del nodo 'A' y tierra (GND). El circuito tiene dos nodos principales, 'A' y 'B'.

\begin{center}
    \begin{circuitikz}[scale=1.2, transform shape]
        \draw (0,0) node[ground] (gnd) {};
        \draw (-4,3) node[circ, label=above:A] (A) {};
        \draw (4,3) node[circ, label=above:B] (B) {};
        
        % Ramas conectadas al nodo A
        \draw (A) to[R, l=$R_2$] (-6,3) to[V, v<=$E_2$] (-6,0) -- (gnd);
        \draw (A) to[R, l=$R_3$] (-4,0) -- (gnd);
        
        % Conexión entre A y B
        \draw (A) to[R, l=$R_5$] (B);
        
        % Ramas conectadas al nodo B
        \draw (B) to[R, l=$R_4$] (4,0) -- (gnd);
        \draw (B) to[R, l=$R_1$] (6,3) to[V, v<=$E_1$] (6,0) -- (gnd);
    \end{circuitikz}
\end{center}

\section{Cálculo de la Tensión de Thévenin ($E_{Th}$)}

La tensión de Thévenin es la tensión en circuito abierto entre A y GND, es decir, $E_{Th} = V_A$. Para calcularla, primero debemos encontrar la tensión en el nodo B, $V_B$.

\subsection{Paso 1: Calcular la Tensión en el Nodo B ($V_B$)}

El nodo B tiene dos ramas en paralelo conectadas a GND: la rama con $R_4$ y la rama con la serie de $E_1$ y $R_1$. Usando el Teorema de Millman en el nodo B:
\begin{equation}
    V_B = \frac{\frac{E_1}{R_1} + \frac{0}{R_4}}{\frac{1}{R_1} + \frac{1}{R_4}} = \frac{\frac{E_1}{R_1}}{\frac{R_1+R_4}{R_1 R_4}} = E_1 \frac{R_4}{R_1 + R_4}
\end{equation}
Esto es un simple divisor de tensión, ya que la rama que va hacia el nodo A está en circuito abierto y no consume corriente.

\subsection{Paso 2: Calcular la Tensión en el Nodo A ($V_A$)}

Ahora que conocemos $V_B$, podemos calcular $V_A$. El nodo A tiene tres ramas en paralelo: la rama con $E_2$ y $R_2$, la rama con $R_3$, y la rama con $R_5$ que se conecta a un punto con potencial $V_B$. Aplicando nuevamente el Teorema de Millman en el nodo A:
\begin{equation}
    V_A = \frac{\frac{E_2}{R_2} + \frac{0}{R_3} + \frac{V_B}{R_5}}{\frac{1}{R_2} + \frac{1}{R_3} + \frac{1}{R_5}}
\end{equation}
Sustituyendo la expresión de $V_B$ que encontramos, obtenemos la fórmula para $E_{Th}$:
\begin{equation}
    E_{Th} = V_A = \frac{\frac{E_2}{R_2} + \frac{1}{R_5}\left(E_1 \frac{R_4}{R_1 + R_4}\right)}{\frac{1}{R_2} + \frac{1}{R_3} + \frac{1}{R_5}}
\end{equation}

\section{Cálculo de la Resistencia de Thévenin ($R_{Th}$)}

Para encontrar $R_{Th}$, apagamos todas las fuentes de tensión (las reemplazamos por cortocircuitos) y calculamos la resistencia equivalente vista desde los terminales A y GND.

\begin{center}
    \begin{circuitikz}[scale=1.2, transform shape]
        \draw (0,0) node[ground] (gnd) {};
        \draw (-4,3) node[circ, label=above:A] (A) {};
        \draw (4,3) node[circ, label=above:B] (B) {};
        
        % Ramas en A
        \draw (A) to[R, l=$R_2$] (-6,3) to[short] (-6,0) -- (gnd);
        \draw (A) to[R, l=$R_3$] (-4,0) -- (gnd);
        \draw (A) to[R, l=$R_5$] (B);
        
        % Ramas en B
        \draw (B) to[R, l=$R_4$] (4,0) -- (gnd);
        \draw (B) to[R, l=$R_1$] (6,3) to[short] (6,0) -- (gnd);
    \end{circuitikz}
\end{center}

Desde el punto de vista del nodo A, tenemos tres caminos en paralelo hacia GND:
\begin{enumerate}
    \item La resistencia $R_2$.
    \item La resistencia $R_3$.
    \item La resistencia $R_5$ en serie con la resistencia equivalente en el nodo B.
\end{enumerate}
La resistencia equivalente en el nodo B ($R_{eqB}$) es el paralelo de $R_1$ y $R_4$:
\begin{equation}
    R_{eqB} = R_1 \parallel R_4 = \frac{R_1 R_4}{R_1 + R_4}
\end{equation}
La resistencia total de la tercera rama es $R_5 + R_{eqB}$.

Finalmente, la resistencia de Thévenin $R_{Th}$ es el paralelo de estas tres ramas:
\begin{equation}
    R_{Th} = R_2 \parallel R_3 \parallel (R_5 + R_{eqB}) = \left( \frac{1}{R_2} + \frac{1}{R_3} + \frac{1}{R_5 + \frac{R_1 R_4}{R_1 + R_4}} \right)^{-1}
\end{equation}

\section{Circuito Equivalente de Norton}

El circuito equivalente de Norton se obtiene directamente a partir de los valores de Thévenin.
\begin{itemize}
    \item \textbf{Resistencia de Norton ($R_N$):} Es igual a la resistencia de Thévenin.
    \begin{equation}
        R_N = R_{Th}
    \end{equation}
    \item \textbf{Corriente de Norton ($I_N$):} Se calcula con la Ley de Ohm.
    \begin{equation}
        I_N = \frac{E_{Th}}{R_{Th}}
    \end{equation}
\end{itemize}

\end{document}